\chapter*{R\'esum\'e}
\initial{L}es \'equipes nationales de r\'eponse aux incidents font face \`a un
d\'es\'equilibre structurel : le volume d'alertes cro\^it plus vite que l'effectif
capable de les traiter, et le d\'elai de r\'eaction reste born\'e par la disponibilit\'e
d'un analyste. Les plateformes d'orchestration existantes n'y r\'epondent
qu'imparfaitement, car elles s'arr\^etent \`a la recommandation : la d\'ecision, donc
l'attente, demeure humaine. Le pr\'esent travail examine la question inverse
--- \`a quelles conditions un syst\`eme peut-il agir seul ? --- et y r\'epond par une
contrainte unique : \emph{la r\'eversibilit\'e est la condition de l'autonomie}. La
plateforme con\c cue, CIRTDEFENSE, ex\'ecute sans validation pr\'ealable les seules
actions qu'elle sait annuler, refuse d'agir sur un contexte non document\'e,
applique une politique r\'edig\'ee en langage naturel et compil\'ee \`a l'avance,
surveille l'effet de ses propres gestes et les annule d'elle-m\^eme s'ils d\'egradent
la cible. Face \`a une menace absente de son catalogue, elle ne se tait pas : elle
d\'eduit un confinement des seuls indicateurs observ\'es, engage ce qui est
r\'eversible et soumet \`a l'humain ce qui laisse une trace durable. Chaque geste est
inscrit dans un journal cha\^in\'e par empreintes, seule trace opposable. La
mod\'elisation suit le processus unifi\'e et le formalisme UML~2.5 ; la r\'ealisation
est valid\'ee par XXX tests automatis\'es, dont XX crit\`eres de recette d\'eriv\'es
directement des exigences.

\paragraph{Mots cl\'es : r\'eponse autonome aux incidents, r\'eversibilit\'e, orchestration
de la s\'ecurit\'e, tra\c cabilit\'e, UML.}

\clearpage
